\documentclass[12pt]{article}
\usepackage{fontspec}
\setmainfont{Liberation Serif}
\usepackage{amsmath,amssymb}
\usepackage[margin=2.4cm]{geometry}
\usepackage{booktabs}
\usepackage{hyperref}

\title{\textbf{Уравнения Навье--Стокса: точная редакция\\ универсального геометрического эффекта $b$}}
\author{wild8highlander}
\date{13 сентября 2026~г.}

\begin{document}
\maketitle

\begin{abstract}
Настоящая монография --- точная редакция: все константы заданы замкнутыми формами,
все числа воспроизводимы из открытого кода верификации L1--L5. Универсальный
геометрический эффект
\[
b=\frac{\pi}{4\pi^{2}+2\pi\sqrt3}=\frac{1}{4\pi+2\sqrt3}=0{,}0623811941210282275463\ldots
\]
возникает аналитически из уравнений Кирхгофа и универсален --- появляется на любой
геометрической поверхности. Поворот скорости на угол
$\theta_b=\arcsin(b)=3{,}5765013142837210\ldots^\circ$ вокруг оси вихря --- \emph{не
внешняя сила}: это естественное поведение жидкости в уравнениях Эйлера, где внутреннее
трение слоёв совершает минимальное движение и сдвиг; $\sin\theta_b=b$ --- точно.
Уравнения остаются исходными уравнениями Навье--Стокса.
\end{abstract}

\section{Точная константная цепочка}

Каждая величина задана замкнутой формой; численные значения получены в mpmath
с 50 знаками и независимо в Julia (BigFloat, 120 бит).

\begin{table}[h]
\centering
\begin{tabular}{lll}
\toprule
Величина & Форма & Значение (50 знаков) \\
\midrule
$b$ & $\pi/(4\pi^{2}+2\pi\sqrt3)$ & $0{,}0623811941210282275\ldots$ \\
$b$ & $1/(4\pi+2\sqrt3)$ & то же (разность форм $=0$) \\
$\theta_b$ & $\arcsin(b)$ & $0{,}0624217236361554344\ldots$ рад $=3{,}5765013\ldots^\circ$ \\
$\sin\theta_b-b$ & тождество & $0$ \emph{точно} \\
$\cos\theta_b$ & $\sqrt{1-b^{2}}$ & $0{,}9980523967307701392\ldots$ \\
$b\,\beta_K L_{\min}$ & $\beta_K=5/3$ & $0{,}3013167052374188596\ldots$ \\
$Z_{\mathrm{full}}/Z_{\mathrm{leading}}$ & $e^{b\beta_K L_{\min}}$ & $1{,}3516373443851241821\ldots$ \\
$C_s$ (Лилли) & $\pi^{-1}(9/4)^{-3/4}$ & $0{,}1732659558297058017\ldots$ \\
$\alpha$ (Клейн) & константа монографии & $2{,}2469796037174670610\ldots$ \\
$L_{\min}$ (Клейн) & константа монографии & $2{,}8981494453551726700\ldots$ \\
\bottomrule
\end{tabular}
\end{table}

Замыкание $e^{b\beta_K L_{\min}}=1{,}3516\ldots$ --- это независимая проверка цифр
константы $L_{\min}$: оба задания (из монографии и из формулы) совпадают на всех 42
выписанных знаках. Эквивалентность двух форм $b$ доказывается символьно (sympy
упрощает разность до нуля).

\section{Поворот --- естественное поведение жидкости, а не внешняя сила}

Смещение на угол не является внешней силой. Это естественное поведение жидкости
в уравнениях Эйлера: внутреннее трение слоёв друг о друга совершает минимальное
движение и сдвиг; доля этого сдвига, приходящаяся на поворот вокруг оси вихря,
равна $b$. Слои жидкости трутся друг о друга --- и этот контакт реализуется именно
как минимальный локальный поворот, без подвода энергии извне.

Точная форма поворота на угол $\varphi$ вокруг локальной оси $\hat\omega$:
\begin{equation}
u'=u_{\parallel}+\sqrt{1-\sin^{2}\varphi}\;u_{\perp}+\sin\varphi\;(\hat\omega\times u_{\perp}),
\qquad u_{\parallel}=(u\cdot\hat\omega)\hat\omega,\quad u_{\perp}=u-u_{\parallel}.
\label{eq:rot}
\end{equation}
При $\varphi=\theta_b$: $\sin\varphi=b$, $\cos\varphi=\sqrt{1-b^{2}}$ --- коэффициенты
формы (\ref{eq:rot}) и есть универсальный эффект $b$. Матрица (\ref{eq:rot}) ---
собственная ортогональная матрица Родригеса: $R^{T}R=I$ точно, $\det R=1$
(уровень L2: невязка $4{,}4\cdot10^{-16}$ на $10^{5}$ векторов).

Никакой модификации уравнений не вводится: полю $u'$ после поворота
продолжает жить в исходных уравнениях Навье--Стокса; градиентная часть,
возникающая при точечном повороте, поглощается давлением (проекция Лерэ).

\section{Кирхгофовское происхождение}

Уравнения Кирхгофа для пары вихрей
$(dx/dt,\,dy/dt)=R(-90^\circ)\,\nabla H$ --- гамильтонов поток: дрейф инварианта
$H$ падает как $dt^{4}$ (RK4, уровень L3), измеренный порядок метода $4{,}00$,
фазовый поток изометричен (жёсткое вращение), площадь сохраняется. Динамика
угла $\varphi$ --- чистый поворот без диссипации; доля поворота $b$ ---
универсальный математический эффект, возникающий на любой геометрической
поверхности: 2D, $S^{2}$, $H^{2}$, $T^{2}$, квартика Клейна, $\mathbb{R}^{3}$,
$S^{3}$.

\section{Точные энергетические тождества (2D)}

Для дивергентно-свободного $u$ поворот $u'=\cos\theta_b\,u+b\,(\hat z\times u)$
(эквивалент формы (\ref{eq:rot}) в 2D) даёт \emph{точно}:
\begin{align}
\omega'&=\cos\theta_b\;\omega, & \text{(невязка } 7{,}1\cdot10^{-14}\text{)}\\
|u'|^{2}&=|u|^{2}, & \text{(}0\text{ точно)}\\
\operatorname{div}u'&=-b\,\omega & \text{(уходит в давление).}
\end{align}
Инжекции энергии нет --- на машинном уровне. В 3D-эталоне максимальное
изменение энергии за один поворот $1{,}4\cdot10^{-9}$ (остаток проекции Лерэ).

\section{3D-эталон: зафиксированный протокол}

Тейлор--Грин в $T^{3}$, $N=48$, $\nu=0{,}01$ ($\mathrm{Re}_{TG}=100$), $T=6$,
$dt=0{,}004$; 2/3-правило через 3/2-паддинг; RK4; непрерывный поворот:
угол за шаг $dt\cdot\theta_b$, кумулятивно $\theta_b T\approx 21{,}459^\circ$;
ось --- направление сглаженного вихря ($k\le N/6$) с квадратичным весом
$(|\omega|/\max|\omega|)^{2}$.

Результаты выполненного прогона (JSON сохраняется при каждом запуске):
истинные NSE --- $I_{BKM}=9{,}3560$, $\max\|\omega\|_\infty=2{,}0000$ (точно ---
максимум-принцип); с непрерывным $b$-поворотом --- фактор $0{,}967\times$ по
BKM-интегралу (нейтрально: изменение $+3{,}3\%$) и $1{,}0000\times$ по
$\max\|\omega\|_\infty$; $|\Delta E|\le1{,}4\cdot10^{-9}$ за поворот (инжекции нет).
Существенное снижение вихря даёт точное 2D-тождество
$\omega'=\cos\theta_b\,\omega$ (уровень L4). Заявлять фактор $>1$ в
зафиксированном протоколе нельзя. \emph{Историческое} значение
«$3{,}5\times$» относится к недоразрешённой конфигурации $N=24$ (глава 11
первой редакции) и не воспроизводится в зафиксированном протоколе; в пакете
оно сохранено с пометкой своей конфигурации.

\section{Многоуровневая верификация}

\begin{table}[h]
\centering
\begin{tabular}{llp{7.2cm}}
\toprule
Уровень & Статус & Содержание \\
\midrule
L1 & PASS & константы из замкнутых форм, 50 знаков, замыкание Клейна \\
L2 & PASS & $R^{T}R=I$, $\det R=1$, спектр $\{1,e^{\pm i\theta_b}\}$ \\
L3 & PASS & гамильтоновость, порядок RK4 $=4{,}00$, изометрия \\
L4 & PASS & тождества поворота в 2D, энергетический баланс \\
L5 & PASS & 3D-эталон: BKM, $\max\|\omega\|_\infty$, без инжекции энергии \\
\bottomrule
\end{tabular}
\end{table}

Каждый уровень пишет JSON-результат при каждом запуске; наборы существуют
на Python (numpy/mpmath) и Julia (чистый stdlib, собственный FFT radix-2
с самотестом против DFT).

\section{Связь с задачей Клэя}

Никакого условия «$b=0$» не требуется и никогда не требовалось: $b$ не является
ни модификацией уравнения, ни внешней силой, ни поправкой в уравнении ---
это название универсального математического эффекта. Поворот --- естественное
поведение жидкости, уже присутствующее в уравнениях Эйлера; уравнения остаются
исходными. Вопросы четырёх режимов Клэя (A: 3D гладкость, B: 3D сингулярность,
C: 2D гладкость, D: 2D сингулярность) формулируются в терминах исходных уравнений;
настоящая редакция даёт для них точную константную базу и воспроизводимый
численный контур, но не заявляет полного решения задачи Клэя.

\section{Сводка исправлений относительно первой редакции}

\begin{itemize}
\item три противоречивых угла ($7{,}07^\circ$ / $4{,}50^\circ$ / $3{,}58^\circ$)
заменены одним: $\theta_b=\arcsin(b)=3{,}5765\ldots^\circ$, $\sin\theta_b=b$ --- точно;
\item значение «$b\approx0{,}0785$» (и форма $\theta_b=b\,\pi/2$) исключены;
точное $b=1/(4\pi+2\sqrt3)=0{,}0624\ldots$;
\item удалены формулировки «требуется $b=0$», «$b$ --- модификация уравнения»,
«$C_K=1{,}5$ --- предсказание»: $C_K=1{,}5$ --- калибровка, согласованная с
независимой оценкой $\gamma=0{,}95456$, $C_K(\gamma)=1{,}4786$;
\item поворот определён как естественное поведение жидкости (внутреннее трение
$\to$ минимальное движение и сдвиг), не как внешнее воздействие;
\item все исторические численные результаты сохранены с пометкой их фактической
конфигурации.
\end{itemize}

\end{document}
